% STATUS: DRAFT
\chapter{Weak closure and the double commutant theorem}\label{ch:doublecomm}

\begin{physmotiv}
A \cstar-algebra contains only \emph{continuous} functions of its observables, so it lacks the yes/no observables (spectral projections) and the macroscopic ``classical'' observables of \cref{ch:topologies}. Closing the algebra in the weak topology fixes this. The remarkable theorem of von Neumann says that this \emph{analytic} closure is the same as a purely \emph{algebraic} operation: take the commutant twice. Two things then follow: von Neumann algebras come in dual pairs $\M,\M'$, and a region of spacetime and its causal complement will give such a pair.
\end{physmotiv}

\section{The commutant}

For a set $\mathcal S\subset\BH$, its \emph{commutant} is $\mathcal S'=\{b\in\BH:ab=ba\ \forall a\in\mathcal S\}$. Properties: $\mathcal S'$ is a unital algebra, weakly closed, and $^*$-closed if $\mathcal S=\mathcal S\ad$; $\mathcal S\subset\mathcal S''$; $\mathcal S_1\subset\mathcal S_2\Rightarrow\mathcal S_2'\subset\mathcal S_1'$; and $\mathcal S'''=\mathcal S'$.

Physically $\mathcal S'$ is ``everything that can be measured simultaneously with all of $\mathcal S$ without disturbing it''. It is the algebra of the \emph{complement} of the system.

\begin{example}
(i) $\M=M_n\otimes\id\subset\Bnd{\C^n\otimes\C^k}$: $\M'=\id\otimes M_k$. Subsystem and environment are mutual commutants. (ii) $\M=\BH$: $\M'=\C\id$. (iii) $\M=\Linf(X,\mu)$ acting by multiplication on $L^2(X,\mu)$ (with $\mu$ $\sigma$-finite): $\M'=\M$ (a maximal abelian algebra, a complete set of commuting observables, like the position operators). (iv) GNS representation of a mixed state of $M_n$ (\cref{ch:gns}): $\pi(\A)'$ is right multiplication, the purifying system.
\end{example}

\section{The double commutant theorem}

\begin{theorem}[von Neumann]\label{thm:bicomm}
Let $\M\subset\BH$ be a unital $^*$-algebra. Then the following are equal:
\[
\M''\;=\;\overline{\M}^{\rm weak}\;=\;\overline{\M}^{\rm strong}\;=\;\overline{\M}^{\sigma\text{-weak}}\;=\;\overline{\M}^{\sigma\text{-strong}}.
\]
\end{theorem}

\begin{definition}
A \emph{von Neumann algebra} is a unital $^*$-algebra $\M\subset\BH$ with $\M=\M''$, equivalently one that is closed in the weak topology.
\end{definition}

\begin{proofsketch}
The inclusion $\overline\M^{w}\subset\M''$ is easy: $\M''$ is weakly closed and contains $\M$. The converse is the content. For a vector $\xi$ let $P$ be the projection onto $\overline{\M\xi}$; since $\M\xi$ is $\M$-invariant, $P\in\M'$. If $T\in\M''$ then $T$ commutes with $P$, so $T$ preserves $\overline{\M\xi}$; and since $\xi=\id\xi\in\M\xi$, we get $T\xi\in\overline{\M\xi}$: it can be approximated by $a\xi$ with $a\in\M$. Doing this for finitely many vectors at once (use the amplification $\HH\otimes\C^k$ with the algebra $\M\otimes\id_k$) gives strong density, hence weak closure.
\end{proofsketch}

\begin{whatwrong}
The theorem needs $\M$ to be a $^*$-algebra containing the identity (or, in general, acting non-degenerately). The \emph{choice of Hilbert space} matters: the weak closure of the \emph{same} abstract \cstar-algebra depends on the representation. For the product states $\psi_\theta$ of the spin chain, $\pi_\theta(\A)''=\BH_\theta$; for $\omega=\tau$ it is the hyperfinite II$_1$ factor; for $p\ne1/2$ it is a type III factor.
\end{whatwrong}

\subsection*{Consequences}
\begin{enumerate}
\item \emph{Spectral projections are in the algebra.} If $a\in\M$ is self-adjoint, all spectral projections $\chi_\Omega(a)$ belong to $\M$ (they are strong limits of continuous functions of $a$). Likewise $|a|$, the partial isometry in the polar decomposition, and any Borel function of $a$.
\item \emph{Kaplansky density.} The unit ball of the \cstar-algebra $\A$ is strongly dense in the unit ball of $\A''$. So no ``hidden'' large operators appear in the closure.
\item \emph{Macroscopic observables.} Limits like $\lim M_N$ (\cref{ch:topologies}) belong to $\pi(\A)''$ and are central.
\item \emph{Normal states.} A state of $\M$ is \emph{normal} (continuous in the $\sigma$-weak topology) iff it is a countable-additive probability: $\omega(\sum_ie_i)=\sum_i\omega(e_i)$ for any family of orthogonal projections. Thus normality is exactly Born's rule with countable additivity, and the right notion of ``physical state'' (\cref{ch:traces}).
\item \emph{Abstract characterization (Sakai).} A \cstar-algebra is (isomorphic to) a von Neumann algebra iff it is the dual of a Banach space; this is why von Neumann algebras are also called W$^*$-algebras and are ``measure-theoretic'' (non-commutative $L^\infty$), whereas \cstar-algebras are ``topological'' (non-commutative $C(X)$).
\end{enumerate}

\begin{keyidea}
Von Neumann algebra $=$ $\M''=\M$ $=$ weakly closed unital $^*$-algebra. Always come in pairs $(\M,\M')$. The commutant is the algebra of the ``complement'', and passage to $\M''$ adds exactly the projections and macroscopic observables that a \cstar-algebra lacks.
\end{keyidea}

\section{Dual pairs and locality}

In relativistic QFT, with $\A(R)$ the algebra of a region $R$ and $R'$ its causal complement (points spacelike to all of $R$), locality says $\A(R')\subset\A(R)'$. \emph{Haag duality} is the stronger statement $\A(R')=\A(R)'$ (\cref{ch:haagkastler}), which holds for wedge regions in free field theories. When it holds, the commutant is literally the algebra of the complementary region, the algebraic version of ``system and environment''. Everything in Part IV (modular theory) treats $\M$ and $\M'$ symmetrically; in the Rindler case $\M=\A(R_{\rm right})$ and $\M'=\A(R_{\rm left})$.

\section*{Exercises}
\begin{exercise}
Compute the commutant of $\M=\{a\otimes\id_k:a\in M_n\}\subset M_n\otimes M_k$ and of $\M=\{\mathrm{diag}(a,a,b):a,b\in M_2\}\subset M_6$. Verify $\M''=\M$ in each case. What is $\M\cap\M'$?
\end{exercise}
\begin{exercise}
Show that every von Neumann algebra on a finite-dimensional space is of the form $\bigoplus_kM_{n_k}\otimes\id_{m_k}$, with commutant $\bigoplus_k\id_{n_k}\otimes M_{m_k}$. (Use that every finite-dimensional $^*$-algebra is a sum of matrix blocks.)
\end{exercise}
\begin{exercise}
Let $\A=C[0,1]$ acting on $L^2[0,1]$ by multiplication. Show that $\A''=\Linf[0,1]$ by showing that $\chi_{[0,1/2]}$ is a strong limit of continuous functions $f_n$ with $0\le f_n\le1$. Is the convergence in norm?
\end{exercise}
\begin{exercise}
Show $\M\mapsto\M'$ reverses inclusions, and that $\M'''=\M'$. If $\M=\M_1\bar\otimes\M_2$ acts on $\HH_1\otimes\HH_2$, argue that $\M'=\M_1'\bar\otimes\M_2'$ (Tomita's commutation theorem; you may assume it).
\end{exercise}
